\section*{Summary Sheet}
We reformulate wildlife protection in Etosha National Park as a weighted loss-flow interdiction problem rather than a patrol-coverage problem. The central question is no longer ``How much of the park can be covered?'' but ``Which entry-to-habitat pathways carry the largest expected biological loss, and where can limited teams, sensors, and UAVs most effectively sever that flow?'' To answer that question, we build a three-layer analytic stack: a threat-transport graph that links ingress points to high-value habitat, a lexicographic barrier-design model that first caps the worst residual hotspot debt and then minimizes total residual loss, and a queueing-based reliability model that sizes field teams against response-time commitments.

Under Etosha's geography, that formulation produces a different operating picture from uniform hotspot coverage. The park is best managed as a gateway--hinge--reserve system. Persistent monitoring belongs on a small set of waterhole-linked hinge zones, repeatable patrol templates should connect eastern entry pressure to the southern relay strip, and mobile aerial capacity should remain uncommitted until uncertainty spikes. Evaluated against the same baseline resources, the resulting plan raises weighted protection coverage from 0.64 to 0.87, increases mean detection probability from 0.47 to 0.66, and lifts the protection index from 0.56 to 0.74 while keeping the implied staffing floor near 34 field rangers.

The contribution of this paper is therefore a re-derived decision object. We optimize residual weighted loss flow and service reliability rather than raw patrol presence. That change in objective changes the management conclusion as well: Etosha should not be defended by spreading effort thinly across a broad corridor, but by suppressing a small set of biologically expensive pathways so that no critical hinge remains simultaneously accessible, unobserved, and slow to serve. The same logic transfers to other protected areas by rebuilding ingress nodes, habitat sinks, and station service rates rather than by copying a fixed patrol map.

{\bf Keywords:} wildlife loss-flow interdiction; ingress-habitat network; lexicographic defense design; queueing reliability; protected-area analytics.

\newpage

\section*{Letter to IMMC}
Dear International Mission for Monitoring Conservation (IMMC),

This paper begins with a different operational premise. Large protected areas do not fail because every hectare is equally hard to protect; they fail because a small number of biologically expensive pathways remain open for too long. In Etosha, those pathways are created by the interaction of entry access, road mobility, water-dependent wildlife concentration, and delayed service from ranger stations. That means the real problem is not blanket patrol coverage. It is weighted loss-flow control.

We therefore rebuild the decision architecture from first principles. First, we convert Etosha into a threat-transport graph in which illegal activity enters through a limited set of ingress points and propagates toward high-value habitat. Second, we measure not only where wildlife value is high, but where movement pathways become operational bottlenecks. Third, we choose patrol templates, sentinel devices, and UAV reserve windows by solving a lexicographic barrier-design problem: reduce the worst remaining hotspot debt first, then clean up total residual loss subject to staffing, time, and station limits. Fourth, we size the field force with a queueing reliability model so that response commitments are tied to service probability rather than to average hours alone.

This reformulation changes the strategic answer. The dominant recommendation is not broad corridor saturation, but a three-tier shield. Eastern ingress pressure should be intercepted before it diffuses inland; southern relay cells should receive repeated patrol templates because they connect multiple downstream habitat clusters; and high-uncertainty sectors should be covered by mobile reserve capacity rather than by permanently tied aerial sorties. Under the baseline resource assumptions, this design outperforms uniform deployment and supports a staffing threshold of roughly 34 field rangers.

The policy implication is equally direct. Managers should protect service reliability before expanding hardware. If enough teams are not available to keep delay exceedance under control, more devices do not rescue the system; they merely detect failure faster. Sensor quality and uptime are the second lever, because hinge-zone monitoring meaningfully compresses residual loss flow. Additional devices matter only after the key hinges are already watched. The sequence of priority is therefore: preserve team depth, protect sentinel uptime, maintain repeatable patrol regularity, and only then expand the hardware footprint.

We also recommend a different accountability lens. Parks should not judge anti-poaching performance by patrol kilometers alone. The better metrics are residual hotspot debt, delay exceedance probability in critical sectors, and the fraction of hinge zones that remain simultaneously covered by a patrol template, a sentinel layer, or an on-call reserve. Those metrics tie strategy directly to field failure modes.

Finally, this framework is transferable because it is morphological rather than cosmetic. What changes from park to park are ingress geometry, habitat sinks, and service times. What persists is the logic: map the transport graph, identify the hinges, cap the worst residual debt, and staff the system to a response commitment. For the IMMC, that is the main recommendation. Encourage conservation strategies that block expensive pathways instead of simply coloring more map area as covered.

Sincerely,\\
Team \#\teamnumber

\newpage
